\documentclass[10pt,twocolumn]{article}
\usepackage{fontspec}
\setmainfont{Times New Roman}
\setsansfont{Helvetica Neue}
\setmonofont{Menlo}
\usepackage[letterpaper,top=0.64in,bottom=0.70in,left=0.67in,right=0.67in,columnsep=0.24in]{geometry}
\usepackage{amsmath,amssymb,amsthm,bm,microtype,booktabs,graphicx,xcolor,enumitem}
\usepackage[colorlinks=true,linkcolor=black,citecolor=black,urlcolor=blue!55!black]{hyperref}
\graphicspath{{results/}}
\setlength{\parindent}{1em}\setlength{\parskip}{0pt}
\setlength{\textfloatsep}{7pt plus 2pt minus 2pt}
\setlength{\floatsep}{6pt plus 2pt minus 2pt}
\setlength{\intextsep}{6pt plus 2pt minus 2pt}
\setlist{nosep,leftmargin=*}\renewcommand{\arraystretch}{1.08}
\newtheorem{proposition}{Proposition}[section]
\newtheorem{theorem}[proposition]{Theorem}
\newtheorem{lemma}[proposition]{Lemma}
\theoremstyle{definition}\newtheorem{remark}[proposition]{Remark}
\newcommand{\R}{\mathbb{R}}
\newcommand{\TV}{\operatorname{TV}}
\newcommand{\diver}{\operatorname{div}}
\newcommand{\grad}{\nabla}
\newcommand{\norm}[1]{\left\lVert#1\right\rVert}
\newcommand{\code}[1]{\texttt{#1}}
\newcommand{\PiD}{\Pi}
\newcommand{\Tf}{T_f}

\title{\vspace{-1.55em}\textbf{A Fast Finite-Flow Meyer--Bregman Decomposition}\\[-1pt]
\large Why the scalar hard jump rings, and how to jump the fused state instead}
\author{Joshuah Rainstar\\[-2pt]\small \texttt{joshuah.rainstar@gmail.com}}
\date{August 2026 revision; historical scope through August 21}

\begin{document}
\twocolumn[
\maketitle\vspace{-1.15em}
\begin{abstract}
We study a fixed-cost acceleration of Meyer cartoon--texture decomposition.
The quality reference is a 64-pass, warm-interleaved, optimized C++
Meyer--Bregman recurrence.  An earlier hard jump was much faster, but it left
oscillatory material in the cartoon and generated signed contour halos.  We
show that these are not separate tuning defects.  If $H^K$ is the finite
scalar contraction used by that jump and $P_K=I-H^K$, its exact texture error
relative to the fused reference is
\[
v_H-v_{64}=P_K(u_{64}-s_1)-H^Kv_{64}+q_\mu .
\]
The first term is a zero-mean structural halo; the second is precisely the
uncompleted carrier tail.  The numerical residual of this identity is below
$1.5\times10^{-13}$ on six surgical controls.

The repair jumps the actual six-field Bregman state
$z=(u,w,t_u^x,t_u^y,t_w^x,t_w^y)$.  From the exact residual
$r=\Tf(z)-z$ and a semismooth derivative $A\in\partial\Tf(z)$, a depth-two
Krylov projection evaluates the finite polynomial
$p_m(A)r=(I+A+\cdots+A^{m-1})r$.  Short ordinary passes refresh the nonlinear
disk-projection chart between jumps.  No edge detector, texture classifier,
source-conditioned threshold, training set, or content law is used.  Fixed
19- and 29-pass-cost schedules are closer to fused pass 64 than ordinary
fused pass 19 and 29, respectively, on every surgical control and on
Barbara, Cameraman, and an analytic crossing.  On an 8-thread M4 Mini, the
two schedules take 11.52 and 18.16 ms at $256^2$, versus 34.60 ms for fused
64, yielding 3.00x and 1.91x speedups.  The public C++ implementation,
Python bindings, demos, exact proof arrays, sample comparisons, and complete
timing records accompany the paper.
\end{abstract}
\vspace{0.4em}
\noindent\textbf{Keywords:} Meyer $G$-norm, cartoon--texture decomposition,
Split Bregman, semismooth flow, Krylov polynomial, spectral fusion.
\vspace{0.85em}
]

\section{Introduction}

Cartoon--texture decomposition asks for $f=\bar u+v$, where coherent
geometry belongs to the effective cartoon $\bar u$ and oscillation belongs to
texture $v$.  This is not a frequency split: sharp interfaces can be cartoon,
and a slowly varying carrier can be texture.  Meyer encoded oscillation as
the divergence of a bounded vector field \cite{meyer2001}.  Aujol et al.
developed a BV--$G$ alternating projection method \cite{aujol2005}; Gilles and
Osher gave a Split-Bregman realization based on two ROF subproblems
\cite{gilles2024}.

Our route began as a speed problem.  Transform planning and reuse first made
a fast STFT super-resolution experiment possible.  A subsequent
two-lattice experiment replaced scalar interpolation by transport between
related measurements.  The work then moved to Meyer decomposition: nested
ROF solves were reduced to a warm, fused recurrence; its state was reduced;
divergences and spectra were maintained; the lower-triangular spectral pair
was fused; and one-axis FACR supplied the arbitrary-shape path.  Figure
\ref{fig:lineage} records this path without importing later projects or
terminology.

\begin{figure*}[t]
\centering\includegraphics[width=.985\textwidth]{method_lineage.png}
\caption{Developmental path used in this paper.  The endpoint is the finite
flow of the Meyer state; later and unrelated work is outside scope.}
\label{fig:lineage}
\end{figure*}

Once the optimized 64-pass result was visually accepted as the decomposition
reference, a second question emerged: can the useful trajectory be compiled
into less work?  The first hard jump appeared attractive because it was
nearly one-shot.  Synthetic controls made its failure unmistakable: edge
halos and texture retained in cartoon.  Attempts to repair the output using
content-dependent admission rules were rejected on principle.  The present
paper instead (i) proves the common source of both defects, (ii) derives a
finite jump on the complete fused state, and (iii) validates fixed schedules
against both equal-cost ordinary iteration and the 64-pass reference.

The contribution is not a claim that 29 operator passes equal 64 nonlinear
passes.  It is a measured Pareto result with an explicit error decomposition:
the finite-flow method removes the two structural errors of the scalar jump,
then exposes what remains as Krylov projection, semismooth chart, and finite
tail error.

\section{Fused reference recurrence}

Let $X=\R^{H\times W}$ with periodic forward gradient $D:X\to X^2$ and
$D^*$ its adjoint.  The discrete isotropic variation is
\begin{equation}
J(u)=\TV(u)=\sum_x \norm{(Du)(x)}_2 .
\end{equation}
Meyer's discrete oscillation norm is
\begin{equation}
\norm v_G=\inf_{v=D^*p}\norm p_{\infty,2},\qquad
\norm p_{\infty,2}=\max_x\norm{p(x)}_2 .
\label{eq:gnorm}
\end{equation}
The associated BV--$G$ formulations and alternating projections are
developed in \cite{meyer2001,aujol2005,gillesmeyer2024}.  We use the same
$[0,255]$ intensity convention and fixed parameters $\lambda=.05$, $\mu=40$.

The optimized reference retains two primal fields and four accumulated
gradient fields:
\begin{equation}
z=(u,w,t_u^x,t_u^y,t_w^x,t_w^y).
\label{eq:state}
\end{equation}
Here $w=f-u-v$ is the finite-pass survivor.  Consequently the public
two-product decomposition is
\begin{equation}
v=f-u-w,\qquad \bar u=f-v=u+w.
\label{eq:effective}
\end{equation}
This convention matters in comparisons: raw $u$ from a finite model is not
the effective cartoon $f-v$.

Define the pointwise disk projection and reflected projection
\begin{equation}
\PiD_a(t)=\begin{cases}t,&|t|\le a,\\a,t/|t|,&|t|>a,
\end{cases}\qquad R_a=I-2\PiD_a,
\label{eq:projection}
\end{equation}
and the screened inverse $S(c,\eta)=(cI+\eta D^*D)^{-1}$.  With
\begin{align}
c_u&=\lambda,&\eta_u&=2\lambda,&a_u&=\eta_u^{-1},\nonumber\\
c_w&=\mu^{-1},&\eta_w&=10\mu^{-1},&a_w&=\eta_w^{-1},
\end{align}
one exact warm fused pass $\Tf:z\mapsto z^+$ is
\begin{align}
u^+&=S(c_u,\eta_u)
 \{c_u(u+w)-\eta_uD^*R_{a_u}(t_u)\},\label{eq:upass}\\
w^+&=S(c_w,\eta_w)
 \{c_w(f-u^+)-\eta_wD^*R_{a_w}(t_w)\},\label{eq:wpass}\\
t_u^+&=Du^++\PiD_{a_u}(t_u),\qquad
t_w^+=Dw^++\PiD_{a_w}(t_w).\label{eq:tpass}
\end{align}

\begin{proposition}[Lower-triangular spectral pass]
For periodic boundaries, equations \eqref{eq:upass}--\eqref{eq:wpass} form a
lower-triangular two-by-two system at every Fourier coefficient.  Both
reflected divergences may be transformed before either inverse transform,
and the two primal spectra may be solved in one range traversal.
\end{proposition}
\begin{proof}
$D^*D$ diagonalizes under the periodic DFT.  If $\sigma_u$ and $\sigma_w$
are the two screened multipliers and $\widehat d_u,\widehat d_w$ the reflected
divergences, then
\begin{align*}
\widehat u^+&=\sigma_u\{c_u(\widehat u+\widehat w)
-\eta_u\widehat d_u\},\\
\widehat w^+&=\sigma_w\{c_w(\widehat f-\widehat u^+)
-\eta_w\widehat d_w\}.
\end{align*}
The second equation depends on the newly computed first variable and no
reverse coupling remains.  This is exactly the paired spectral loop in the
C++ kernel.
\end{proof}

The first pass starts from zero state and has a still simpler source-only
triangle.  Later passes reuse primal spectra and four reflected-dual fields.
This warm interleaving differs computationally from cold nested ROF solves,
but it retains the same Meyer--Bregman geometry and provides the high-quality
finite trajectory studied here.

\section{Exact failure mechanism of the hard jump}

The rejected hard construction replaced the state recurrence by two scalar
observations.  Let $\mathcal S$ be its scalar separator, $H$ its contraction,
$P_K=I-H^K$, and
\begin{align}
s_0&=\mathcal S(f),&m_0&=P_K(f-s_0),\nonumber\\
r_1&=f-m_0,&s_1&=\mathcal S(r_1).
\end{align}
Before its final $G$-capacity correction, the texture proposal is
$m=P_K(f-s_1)$.  Write $v_H=m+q_\mu$, where $q_\mu$ is the change caused by
that final capacity route.

\begin{theorem}[Common cause of halo and retained texture]
Let $(u_*,v_*)$ be any exact two-product reference, $f=u_*+v_*$.  The second
hard observation and its final texture error satisfy
\begin{align}
r_1&=u_*+H^Kv_*+P_K(s_0-u_*),\label{eq:secondobs}\\
v_H-v_*&=P_K(u_*-s_1)-H^Kv_*+q_\mu .\label{eq:defect}
\end{align}
\end{theorem}
\begin{proof}
Since $P_K=I-H^K$,
\begin{align*}
r_1&=f-P_K(f-s_0)=s_0+H^K(f-s_0)\\
&=u_*+H^Kv_*+P_K(s_0-u_*),
\end{align*}
which proves \eqref{eq:secondobs}.  Similarly,
\begin{align*}
m-v_*&=P_K(f-s_1)-v_*\\
&=P_K(u_*-s_1)+(P_K-I)v_*\\
&=P_K(u_*-s_1)-H^Kv_*.
\end{align*}
Adding $q_\mu=v_H-m$ proves \eqref{eq:defect}.
\end{proof}

The two visible failures are therefore one representation error.
$H^Kv_*$ is the uncompleted carrier tail; it remains in cartoon with the
opposite sign in texture.  If $e=u_*-s_1$ contains a step mismatch and
$H^Ke=h_K*e$ for a normalized smoothing kernel $h_K$, then
$P_Ke=e-h_K*e$ has the paired positive and negative contour response of a
halo.  It has zero mean, so a one-sided clamp cannot remove it faithfully.
For a Fourier carrier with contraction gain $0<\rho<1$,
$H^Kv_*=\rho^Kv_*\ne0$ at every finite $K$.

Figure \ref{fig:defect} renders the terms independently.  Across six
$256^2$ controls, direct substitution into \eqref{eq:secondobs} and
\eqref{eq:defect} leaves at most $1.42\times10^{-13}$ in $L_\infty$.  This is
an identity check, not a fitted explanation.

\begin{figure*}[p]
\centering\includegraphics[width=.99\textwidth]{hard_jump_defect_atlas.png}
\caption{Surgical decomposition of the old hard-jump error.  The last column
is the sum of the retained-carrier, structural-halo, and capacity-correction
columns and agrees with hard texture minus fused-64 texture to roundoff.
Pure edge and pure carrier isolate the two terms; compound controls show their
superposition.}
\label{fig:defect}
\end{figure*}

This theorem also rules out the discarded content gate as a fundamental
repair.  A classifier applied after $m$ is formed can hide samples of the
error, but it cannot restore the omitted six-field trajectory and must encode
assumptions about image content.  We instead change the representation being
jumped.

\section{Finite flow of the complete state}

After a short ordinary prefix, fix the current state $z$ and define
\begin{equation}
r=\Tf(z)-z,\qquad A\in\partial\Tf(z).
\label{eq:residual}
\end{equation}
If $\Tf$ were affine on the visited chart, its exact finite displacement after
$m$ more passes would be
\begin{equation}
\Tf^m(z)-z=p_m(A)r,\qquad
p_m(A)=I+A+\cdots+A^{m-1}.
\label{eq:finitepoly}
\end{equation}
This is the correct target.  The stationary Newton displacement
$(I-A)^{-1}r$ was tested and rejected: it can overshoot a desired finite
state at structural transitions.

The only nonlinear branch derivative comes from the disk projection already
present in \eqref{eq:projection}.  For $n=t/|t|$,
\begin{equation}
D\PiD_a(t)h=\begin{cases}
h,&|t|<a,\\[2pt]
\dfrac{a}{|t|}(I-nn^T)h,&|t|>a,
\end{cases}
\label{eq:diskderivative}
\end{equation}
and $DR_a(t)h=h-2D\PiD_a(t)h$.  At $|t|=a$ the implementation chooses the
interior member of the Clarke derivative.  Differentiating the fused triangle
gives
\begin{align}
\dot u^+&=S_u\{c_u(\dot u+\dot w)
-\eta_uD^*DR_{a_u}(t_u)\dot t_u\},\nonumber\\
\dot w^+&=S_w\{-c_w\dot u^+
-\eta_wD^*DR_{a_w}(t_w)\dot t_w\},\label{eq:tangent}\\
\dot t_u^+&=D\dot u^++D\PiD_{a_u}(t_u)\dot t_u,\nonumber\\
\dot t_w^+&=D\dot w^++D\PiD_{a_w}(t_w)\dot t_w.\nonumber
\end{align}
No statistic of $f$ enters \eqref{eq:tangent}; only the current optimization
state and the original projection geometry do.

\subsection{Depth-two nonorthogonal Krylov form}

The implementation evaluates $r$, $Ar$, and $A^2r$.  Let
\begin{equation}
\begin{gathered}
B=[r,Ar],\qquad G=B^TB,\\
K=B^TAB,\qquad \widehat A=G^{-1}K.
\end{gathered}
\label{eq:galerkin}
\end{equation}
Because $Ar$ is the second basis vector, the first column of $\widehat A$ is
exactly $(0,1)^T$.  Only the projection of $A^2r$ needs a two-by-two solve.
The applied displacement is
\begin{equation}
\delta_m=Bp_m(\widehat A)e_1.
\label{eq:krylovjump}
\end{equation}
This is algebraically the depth-two Arnoldi polynomial without normalization
or full-field modified Gram--Schmidt sweeps.  A fixed 64-chunk reduction
computes the five required inner products in a thread-count-independent
summation order.

\subsection{What error remains}

Write the semismooth expansion at $z$ as
\begin{equation}
\Tf(z+d)=\Tf(z)+Ad+N_z(d).
\end{equation}
Let $d_k=\Tf^k(z)-z$, $p_0=0$, $p_{k+1}=r+Ap_k$, and
$e_k=d_k-p_k$.  Then
\begin{equation}
e_{k+1}=Ae_k+N_z(d_k).
\end{equation}
If $\norm A\le\gamma$ and $\norm{N_z(d)}\le L\norm d^2/2$ on the traversed
chart, induction gives
\begin{equation}
\norm{e_m}\le \frac{L}{2}\sum_{k=0}^{m-1}
\gamma^{m-1-k}\norm{d_k}^2.
\label{eq:remainderbound}
\end{equation}
The implemented error additionally contains the depth-two projection term
\begin{equation}
\norm{p_m(A)r-Bp_m(\widehat A)e_1}.
\end{equation}
Ordinary passes after each jump refresh the nonlinear disk branches and
shorten the chart over which \eqref{eq:remainderbound} accumulates.  Thus the
remaining discrepancy from pass 64 consists of finite target tail, Krylov
projection, and semismooth chart error.  The terms $P_K(u_*-s_1)$ and
$-H^Kv_*$ of \eqref{eq:defect} are structurally absent: this method never
forms $s_1$, $P_K$, or a scalar surrogate for $z$.

\subsection{Fixed schedules and the shorter-jump ablation}

For prefix $p$, two tangent products, $s$ ordinary refresh passes, and $j$
jumps, the exact operator-pass cost is
\begin{equation}
C=p+j(1+2+s),
\label{eq:cost}
\end{equation}
where the one counts the exact nonlinear residual evaluation.  The public
schedules are fixed independently of the source:
\begin{center}\small
\begin{tabular}{lrrrrr}\toprule
&$p$& horizon &$s$&$j$&$C$\\\midrule
fast&4&18&2&3&19\\
quality&4&10&2&5&29\\\bottomrule
\end{tabular}
\end{center}

We also tested the natural alternative ``jump less, refresh once.''  A cheap
post-jump graph retraction,
\begin{equation}
t\leftarrow Du+\PiD_a(t-Du),
\label{eq:retract}
\end{equation}
restores the dual feasibility $|t-Du|\le a$ in $O(HW)$ pointwise work and no
FFT.  Retraction alone prevents instability but loses most of the
acceleration.  Shorter jumps with one exact refresh won at $128^2$ but did
not uniformly beat the two-refresh schedules at $256^2$.  Reducing only the
horizon traded small wins among controls rather than producing dominance.
This negative ablation shows that refreshes are nonlinear chart updates, not
backtracking overhead.  A useful constrained variant is therefore likely to
cost about the same total spectral work unless it exploits a stronger bound
than \eqref{eq:retract}.

\section{C++ realization}

The public full-spectral entry point is
\path{bfft_meyer_split_flow_jump}.  Its plan owns the ordinary reduced
state and lazily allocates three six-field Krylov vectors plus primal spectra.
The implementation performs:
\begin{enumerate}[label=\arabic*.]
\item a source FFT and ordinary fused prefix;
\item one exact residual pass without overwriting live spectra;
\item two streamed tangent passes using \eqref{eq:diskderivative};
\item one fixed-chunk five-scalar Gram reduction;
\item the $2\times2$ polynomial in registers and one state AXPY;
\item fixed ordinary refresh passes, repeated by schedule.
\end{enumerate}
No candidate scan or line search is performed.  Texture is emitted as
$f-u-w$ and cartoon as $f-v$, guaranteeing exact two-product recomposition.

The Python binding exposes \path{MeyerPlan.split_flow_jump}; the default
full-spectral \code{split} uses the quality schedule.  The effect API and
still-image demo expose \code{quality}, \code{fast}, \code{fused64}, and the
old hard control.  The video demo uses the fixed fast schedule.  FACR plans
retain their configured ordinary fused recurrence because the current
finite-flow tangent is full-spectral; they also fold the finite survivor into
effective cartoon.

Native tests compare the C++ output to an independent NumPy implementation of
\eqref{eq:upass}--\eqref{eq:krylovjump}, check exact recomposition, translation
equivariance, and bit identity across thread counts, and require both public
schedules to beat their equal-cost ordinary controls on the surgical battery.

\section{Evidence}

\subsection{Protocol}

All methods use $\lambda=.05$, $\mu=40$ and are applied unchanged to every
scene.  Fused 64 is the accepted finite-trajectory reference.  The error
reported below is
\begin{equation}
E_{64}(v)=\norm{v-v_{64}}_2/\norm{v_{64}}_2.
\label{eq:error}
\end{equation}
For the cold nested Gilles record, its model survivor $f-u-v$ is folded into
effective cartoon $f-v$ before display, so every column is an exact
two-product pair.  Native timings are medians after three warmups and eleven
repeats on an 8-thread Apple M4 Mini.  Plan construction is excluded.  The
cold Gilles values are archived NumPy/SciPy wall times, capped at 120 outer
iterations; they are historical context, not a same-implementation speed
claim.

\subsection{Surgical controls}

At $256^2$, the old hard jump has mean texture RMS error 3.431 to fused 64
over pure edge, pure carrier, symmetric support, checker support, multiscale
crossing, and junction texture.  The fast flow lowers that to 1.400, versus
1.496 for ordinary fused 19.  The quality flow lowers it to 1.039, versus
1.196 for ordinary fused 29.  The flow/equal-cost ratios lie in
$[.919,.951]$ for fast and $[.821,.927]$ for quality; every scene improves.
Figure \ref{fig:surgical} shows that the large signed contour pattern and
carrier deficit of the old hard output are replaced by the much smaller
ordinary finite-convergence discrepancy.

\begin{figure*}[p]
\centering\includegraphics[width=.99\textwidth]{finite_flow_surgical_atlas.png}
\caption{Fused-64 reference, old hard jump, and 29-cost finite-flow quality
schedule on six surgical controls.  The final two columns use the same
diverging scale within each row.  The flow discrepancy is faint and follows
ordinary finite convergence; the hard discrepancy reproduces the two proven
terms of Figure \ref{fig:defect}.}
\label{fig:surgical}
\end{figure*}

\subsection{Natural and historical comparisons}

Figure \ref{fig:natural} compares effective cartoons and textures.  The hard
control visibly retains Barbara's cloth and Cameraman's ground texture in
cartoon, while its synthetic contours ring.  Both finite-flow outputs track
the fused-64 allocation.  Table \ref{tab:natural} reports the corresponding
times and errors.  At equal operator cost, finite flow is closer than ordinary
fused iteration for all three sources and both schedules.

\begin{figure*}[p]
\centering\includegraphics[width=.985\textwidth]{natural_flow_comparison.png}
\caption{Exact two-product comparison.  Each method is displayed as effective
cartoon $f-v$ and texture $v$.  The cold nested Gilles output is retained as
historical context; fused pass 64 is the trajectory reference.  Texture rows
share a symmetric scale within each source.}
\label{fig:natural}
\end{figure*}

\begin{table*}[t]
\centering\scriptsize
\caption{M4 native median time in ms and texture error $E_{64}$.  Gilles
times (marked $\dagger$) are archived cold NumPy/SciPy runs and are not native
C++ timings.}
\label{tab:natural}
\begin{tabular}{lrrrrrr}\toprule
&\multicolumn{2}{c}{Barbara $512^2$}&\multicolumn{2}{c}{Cameraman $256^2$}
&\multicolumn{2}{c}{analytic $256^2$}\\
method & ms &$E_{64}$& ms &$E_{64}$& ms &$E_{64}$\\\midrule
old hard control &7.418&.5235&2.821&.6923&2.852&.4200\\
ordinary fused 19&25.583&.1080&10.347&.1863&10.182&.1889\\
finite flow fast&33.933&.0923&11.959&.1582&11.869&.1689\\
ordinary fused 29&39.545&.0710&15.791&.1273&15.849&.1385\\
finite flow quality&53.083&.0513&18.144&.0917&18.301&.1111\\
ordinary fused 64&86.988&0&34.864&0&34.732&0\\
cold nested Gilles$^\dagger$&169582.8&.0744&46167.9&.2095&45659.2&.1584\\
\bottomrule
\end{tabular}
\end{table*}

The flow schedules cost more wall time than their equal-pass ordinary
controls because tangent projection and deterministic Gram reduction have
higher constants than an ordinary fused pass.  They nevertheless move
farther along the quality frontier.  At $256^2$, quality reaches $E_{64}=.092$
on Cameraman in 18.14 ms, whereas ordinary 29 reaches $.127$ in 15.79 ms and
fused 64 reaches zero by definition in 34.86 ms.  The choice is an explicit
Pareto trade, not an assertion that pass-equivalent cost equals elapsed time.

\begin{figure*}[t]
\centering\includegraphics[width=.98\textwidth]{finite_flow_pareto.png}
\caption{Time--quality plane.  Blue diamonds are the two proposed fixed
schedules; gray circles are ordinary fused 19, 29, and 64.  The hard jump is
cheap but far from the accepted decomposition.  Cold nested Gilles is shown
only to preserve the historical comparison.}
\label{fig:pareto}
\end{figure*}

\subsection{Scaling}

Table \ref{tab:scaling} and Figure \ref{fig:scaling} use a dependency-free
multiscale crossing at three sizes.  The quality schedule is 2.01x, 1.91x,
and 1.64x faster than fused 64; the fast schedule is 3.05x, 3.00x, and 2.55x
faster.  Absolute 512-square times are 33.92 ms (fast), 52.63 ms (quality),
and 86.58 ms (fused 64).

\begin{table}[h]
\centering\scriptsize
\caption{M4 medians (ms), 8 threads.}
\label{tab:scaling}
\begin{tabular}{rrrrr}\toprule
size&fast&quality&fused 64&speedups\\\midrule
128&6.671&10.107&20.325&3.05/2.01\\
256&11.521&18.155&34.598&3.00/1.91\\
512&33.920&52.633&86.581&2.55/1.64\\
\bottomrule\end{tabular}
\end{table}

\begin{figure}[h]
\centering\includegraphics[width=.98\columnwidth]{native_flow_scaling.png}
\caption{Native scaling, including equal-cost ordinary controls.}
\label{fig:scaling}
\end{figure}

\section{Scope and limitations}

Fused pass 64 is a finite, visually accepted reference, not a theorem of
ground-truth uniqueness.  The finite-flow schedules approximate its
trajectory; they do not solve every BV--$G$ formulation or equal the cold
nested Gilles state.  The evidence covers one parameter pair, periodic
full-spectral power-of-two shapes, six designed controls, and three comparison
images.  FACR and Neumann paths remain ordinary fused methods.

The depth-two derivative freezes one semismooth chart.  Projection-boundary
crossings generate the remainder in \eqref{eq:remainderbound}; refresh passes
limit but do not eliminate it.  Larger Krylov spaces, certified chart radii,
or a useful low-cost retraction could improve the frontier, but the shorter-
jump ablation warns against selecting a schedule from one resolution.  The
old hard control remains public only for reproduction and negative tests.

\section{Reproduction}

The proof is generated by
\path{experiments/meyer_hard_jump_defect_proof.py}; the independent
finite-flow oracle and schedule battery are in
\path{experiments/meyer_semismooth_state_jump.py}.  Native scaling is
serialized by \path{experiments/benchmark_meyer_flow_jump.py}.  The paper
directory contains the archived Gilles arrays, the current natural benchmark,
all timing samples, and a renderer that reads records rather than recomputing
claims.  The ABI, implementation, binding, demo paths, and invariance tests are
in the main source tree.

\section{Conclusion}

The old hard jump failed for a precise reason: it compressed a six-field
nonlinear trajectory into two finite scalar observations.  Its halo and
texture leakage are the two exact terms of that replacement.  Once the jump
is applied to the actual fused state, those terms disappear from the error
algebra.  A finite Krylov polynomial then advances the same Meyer--Bregman
recurrence, while ordinary refreshes control the nonlinear chart remainder.

The result is a fixed, content-free, reproducible Pareto improvement.  It
does not obtain speed by deciding what an image contains.  It obtains speed
by representing more of the work that the accepted decomposition is already
doing.

\section*{AI-assisted research disclosure}

Joshuah Rainstar supplied the research direction, visual judgments, and
successive falsification criteria.  AI systems generated substantial code,
algebraic exploration, benchmark scaffolding, and manuscript text under that
direction.  Claims are tied to executable artifacts and serialized records.

\begin{thebibliography}{9}
\bibitem{meyer2001} Y. Meyer, \emph{Oscillating Patterns in Image Processing
and in Some Nonlinear Evolution Equations}. AMS, 2001.
\bibitem{aujol2005} J.-F. Aujol, G. Aubert, L. Blanc-F\'eraud, and A.
Chambolle, ``Image decomposition into a bounded variation component and an
oscillating component,'' \emph{J. Math. Imaging Vision}, 22:71--88, 2005.
\bibitem{chambolle2004} A. Chambolle, ``An algorithm for total variation
minimization and applications,'' \emph{J. Math. Imaging Vision}, 20:89--97,
2004.
\bibitem{goldstein2009} T. Goldstein and S. Osher, ``The Split Bregman method
for $L^1$-regularized problems,'' \emph{SIAM J. Imaging Sci.}, 2:323--343,
2009.
\bibitem{gilles2024} J. Gilles and S. Osher, ``Bregman implementation of
Meyer's $G$-norm for cartoon + textures decomposition,'' arXiv:2410.22777,
2024.
\bibitem{gillesmeyer2024} J. Gilles and Y. Meyer, ``Properties of BV--G
decomposition models,'' arXiv:2411.04456, 2024.
\end{thebibliography}
\end{document}
